\documentclass[11pt]{article}
\usepackage[margin=1in]{geometry}
\usepackage[T1]{fontenc}
\usepackage[utf8]{inputenc}
\usepackage{lmodern}
\usepackage{microtype}
\usepackage{array}
\usepackage{longtable}
\usepackage{booktabs}
\usepackage{enumitem}
\usepackage{xcolor}
\usepackage{hyperref}

\hypersetup{
    colorlinks=true,
    linkcolor=blue,
    urlcolor=blue,
    citecolor=blue
}

\setlength{\parindent}{0pt}
\setlength{\parskip}{0.7em}
\renewcommand{\arraystretch}{1.18}
\setlist[itemize]{leftmargin=1.5em}
\setlist[enumerate]{leftmargin=1.8em}

\newcommand{\coursecode}{MATH 103.1}
\newcommand{\coursetitle}{Predictive Modeling for Text}
\newcommand{\policytitle}{GenAI Policy for \coursecode}

\begin{document}

\begin{center}
    {\Large\bfseries \policytitle}\par
    \vspace{0.25em}
    {\large \textit{\coursetitle}}\par
    \vspace{0.5em}
    {Version 1.0}
\end{center}

\section*{Policy Statement}

Generative AI tools such as ChatGPT, GitHub Copilot, Gemini, Claude, and similar systems may be used in \coursecode: \coursetitle{} as learning, coding, debugging, documentation, research-support, and critique assistants.

The use of GenAI is not prohibited. However, its use must be visible, bounded, auditable, and pedagogically meaningful.

This policy is guided by the principle that GenAI should support human learning, not replace it. The key question is not merely how students can use AI, but when, where, and how AI can meaningfully contribute to human learning.

In this course, GenAI use is pedagogically meaningful when it helps students better understand, implement, evaluate, interpret, and critique text prediction experiments without replacing the student's responsibility as analyst, experimenter, and ethical decision-maker.

\section{Guiding Principle}

\textbf{GenAI may assist the experiment, but it may not replace the experimenter.}

\coursecode{} treats text prediction as an experimental process. Students must learn how to frame a problem, prepare text data, choose representations, train models, evaluate results, interpret errors, and defend conclusions. GenAI may help students during this process, but the student remains responsible for the final work.

\section{Acceptable Uses of GenAI}

Students may use GenAI tools for the following purposes, provided that the use is disclosed when it materially affects submitted work and that the student verifies the output before relying on it.

\begin{itemize}
    \item \textbf{Concept clarification.} Students may ask GenAI to explain concepts such as tokenization, TF-IDF, embeddings, logistic regression, support vector machines, neural networks, macro-F1, cross-validation, overfitting, data leakage, model interpretation, or error analysis.
    \item \textbf{Coding support.} Students may ask GenAI to suggest Python code, explain syntax, debug errors, improve readability, or help organize Jupyter notebook cells.
    \item \textbf{Workflow critique.} Students may ask GenAI to critique a modeling plan, identify possible data leakage, suggest evaluation metrics, or compare alternative preprocessing choices.
    \item \textbf{Documentation support.} Students may use GenAI to improve comments, docstrings, markdown explanations, and report organization, provided that the final explanation reflects the student's own understanding.
    \item \textbf{Reflection support.} Students may use GenAI to help formulate reflection questions or identify possible limitations of their work, but the actual reflection must be based on the student's own experience and judgment.
    \item \textbf{Source and documentation navigation.} Students may ask GenAI to help locate relevant concepts in course notes, textbooks, or official documentation, but must independently verify technical claims, citations, quotations, or numerical information.
\end{itemize}

\section{Risk Levels for GenAI Use}

To help students distinguish responsible assistance from inappropriate substitution, GenAI use in this course is classified into three zones.

\begin{longtable}{p{0.18\textwidth}p{0.32\textwidth}p{0.42\textwidth}}
\toprule
\textbf{Zone} & \textbf{Status} & \textbf{Examples} \\
\midrule
\endfirsthead
\toprule
\textbf{Zone} & \textbf{Status} & \textbf{Examples} \\
\midrule
\endhead
\bottomrule
\endfoot
Green zone & Allowed; disclose when used in submitted work if it affected the submission. & Concept clarification, syntax explanation, grammar polishing, debugging guidance, formatting suggestions. \\
Yellow zone & Allowed only with verification, revision, and documentation. & AI-suggested code blocks, modeling workflows, feature engineering choices, metric selection, interpretation of results, literature or source summaries. \\
Red zone & Not allowed. & Fabricated results, fake citations, hidden AI-written reports, AI completion of an entire lab/project/exam answer, undisclosed AI-generated work that the student cannot explain. \\
\end{longtable}

\section{Unacceptable Uses of GenAI}

The following uses are not allowed:

\begin{itemize}
    \item Submitting AI-generated code, explanations, or reports without understanding, testing, and revising them.
    \item Asking GenAI to complete an entire laboratory activity, project report, critique submission, or examination answer on the student's behalf.
    \item Copying GenAI-generated interpretations of model results without checking whether they are supported by the actual output.
    \item Using GenAI to fabricate results, citations, experiment logs, GitHub history, MLflow runs, reflections, or data descriptions.
    \item Using GenAI to hide lack of participation in group work.
    \item Submitting work where the student cannot explain the code, modeling choices, evaluation results, or conclusions.
    \item Uploading confidential, private, sensitive, copyrighted, or restricted data to GenAI tools unless the data are publicly available or the instructor has explicitly permitted such use.
\end{itemize}

\section{Examination Rule}

During examinations, GenAI use is prohibited unless the instructor explicitly states otherwise in writing for a specific examination component. Any permitted use will be announced together with the examination instructions and will be governed by the same principles of visibility, verification, and human responsibility.

\section{Required Disclosure}

All GenAI use that materially affects a submitted notebook, critique, report, code file, analysis, interpretation, reflection, or project artifact must be disclosed through a GenAI Use Log.

A student who did not use GenAI should state: ``No GenAI tool was used for this submission.''

Material GenAI use means AI assistance that influenced the submitted code, explanation, analysis, model choice, interpretation, report structure, reflection, or final decision. Minor uses may be summarized in one log entry when several prompts served the same purpose.

\begin{longtable}{p{0.33\textwidth}p{0.59\textwidth}}
\toprule
\textbf{Required Information} & \textbf{Student Response} \\
\midrule
\endfirsthead
\toprule
\textbf{Required Information} & \textbf{Student Response} \\
\midrule
\endhead
\bottomrule
\endfoot
GenAI tool used & ChatGPT, GitHub Copilot, Gemini, Claude, or other tool. \\
Purpose of use & Explanation, debugging, code suggestion, critique, documentation, source navigation, reflection, or other purpose. \\
Prompt or summary of prompt & What the student asked. For repeated minor uses, a concise summary is acceptable. \\
Useful output & What was helpful. \\
Rejected, revised, or corrected output & What was incomplete, wrong, misleading, unsuitable, or changed by the student. \\
Verification step & How the student checked correctness. \\
Human decision & What the student finally decided as analyst and experimenter. \\
\end{longtable}

\section{Verification Requirement}

Students must verify any AI-assisted output before using it. Verification may include:

\begin{itemize}
    \item checking the lecture notes or textbook;
    \item running the code and inspecting outputs;
    \item testing code on a small example;
    \item checking Python, pandas, PyTorch, or scikit-learn documentation;
    \item comparing model results across train, validation, and test sets;
    \item checking whether preprocessing caused data leakage;
    \item checking whether metrics match the prediction task and class distribution;
    \item explaining the result in the student's own words;
    \item asking whether the result makes sense mathematically, statistically, and ethically.
\end{itemize}

A working code cell is not enough. Students must be able to explain why the code is appropriate and how the result should be interpreted.

\section{Data Privacy and Source Integrity}

Students must not upload confidential, private, sensitive, copyrighted, or restricted data to GenAI tools unless the data are publicly available or the instructor has explicitly permitted such use. When in doubt, students should use small synthetic examples, anonymized excerpts, or instructor-approved datasets.

GenAI-generated citations, quotations, dataset descriptions, legal claims, technical claims, or numerical results must be verified against reliable sources such as course notes, textbooks, official documentation, peer-reviewed sources, or instructor-provided materials. Students may not cite a source merely because a GenAI tool suggested it.

\section{Human Judgment Requirement}

Every major laboratory or project submission must show evidence of human judgment. Students must explain:

\begin{enumerate}
    \item the prediction problem being solved;
    \item the target variable and input text;
    \item the preprocessing and vectorization choices;
    \item the model or models used;
    \item the evaluation metric and why it is appropriate;
    \item the main errors or limitations of the model;
    \item how GenAI helped, misled, or was corrected;
    \item what final decisions were made by the student.
\end{enumerate}

This requirement ensures that GenAI use remains pedagogically meaningful.

\section{Application to Laboratory Notebooks}

Each Jupyter notebook must contain a required section titled \textbf{GenAI Use Disclosure}.

Required notebook template:

\begin{verbatim}
## GenAI Use Disclosure

1. Did you use any GenAI tool for this notebook?
   Answer:

2. What tool did you use?
   Answer:

3. What did you use it for?
   Answer:

4. Provide the prompt or a short summary of the prompt.
   Answer:

5. What part of the AI response did you use?
   Answer:

6. What part of the AI response did you reject, revise, or correct?
   Answer:

7. How did you verify correctness?
   Answer:

8. What final decision did you make as the student-researcher?
   Answer:
\end{verbatim}

\section{Application to Weekly Critique Submissions}

For weekly critique submissions, GenAI may be used to clarify readings, generate questions, or improve writing. However, the critique must reflect the student's own analysis.

Students may not submit a generic AI-generated summary of a reading, notebook, or model result. A valid critique should show:

\begin{itemize}
    \item the student's own understanding;
    \item a specific insight, question, or objection;
    \item connection to the week's lecture or lab;
    \item evidence of independent judgment;
    \item disclosure of GenAI use, if any.
\end{itemize}

\section{Application to the Capstone Project}

For the capstone project, GenAI may be used as a coding assistant, debugging assistant, research-support assistant, documentation assistant, or critique partner. However, the final project must include:

\begin{enumerate}
    \item a reproducible notebook or code repository;
    \item clear data documentation;
    \item train-validation-test split or equivalent validation design;
    \item model evaluation;
    \item error analysis;
    \item interpretation of results;
    \item GenAI Use Log;
    \item reflection on Human-AI collaboration.
\end{enumerate}

Students must be prepared to explain their project orally or in writing. The instructor may conduct brief oral or written checks asking students to explain selected code cells, modeling choices, evaluation results, or GenAI-assisted portions of their submission.

\section{Group Work Accountability}

In group submissions, each member remains responsible for understanding the parts of the work they contributed to and for accurately disclosing any GenAI assistance used in those parts. The group GenAI Use Log should identify whether AI-assisted work affected code, analysis, writing, visualization, debugging, project planning, or interpretation.

A group may not use GenAI to mask non-participation. Each member may be asked to explain the parts of the submitted work for which they claim responsibility.

\section{Academic Integrity and Consequences}

Failure to disclose GenAI use that materially affected the submitted code, explanation, analysis, report, interpretation, reflection, or project artifact may be treated as an academic integrity concern.

Minor or first-time disclosure problems may require revision, resubmission, or oral explanation. Serious cases, including fabricated results, fabricated citations, hidden AI-generated submissions, or inability to explain submitted work, may be handled under the university's academic-integrity procedures.

More importantly, submitting AI-generated work without understanding violates the purpose of the course. The course is designed to develop the student's ability to conduct, interpret, and critique machine learning experiments involving text data.

The issue is not whether GenAI was used. The issue is whether the student used GenAI responsibly and whether the final submission demonstrates genuine learning.

\section{Instructor's Evaluation Standard}

Student work will be evaluated not only on whether the code runs or whether the answer appears correct, but also on whether the submission shows the following:

\begin{longtable}{p{0.32\textwidth}p{0.60\textwidth}}
\toprule
\textbf{Criterion} & \textbf{Description} \\
\midrule
\endfirsthead
\toprule
\textbf{Criterion} & \textbf{Description} \\
\midrule
\endhead
\bottomrule
\endfoot
Technical correctness & The code, model, and results are correct and appropriate. \\
Reproducibility & The workflow can be rerun and checked. \\
Conceptual understanding & The student can explain what was done and why. \\
Evaluation quality & Metrics and validation design are appropriate. \\
Error analysis & The student examines mistakes and limitations. \\
Source integrity & Claims, citations, and external information are verified. \\
GenAI transparency & AI use is clearly disclosed. \\
Human judgment & The final decisions are made and justified by the student. \\
\end{longtable}

\section{Closing Statement}

GenAI tools such as ChatGPT, GitHub Copilot, Gemini, Claude, and similar systems may be used in \coursecode{} as learning, coding, debugging, documentation, research-support, and critique assistants. However, GenAI use must be visible, bounded, auditable, and pedagogically meaningful.

Students must disclose material GenAI use through a GenAI Use Log and must verify any AI-assisted output before including it in a submission. Students remain fully responsible for the correctness, reproducibility, interpretation, and ethical defensibility of their work. GenAI may assist the experiment, but it may not replace the experimenter.

In this course, responsible GenAI use is part of learning modern predictive analytics. Students are expected to use AI tools critically, transparently, and ethically while developing their own competence as analysts, experimenters, and interpreters of text prediction models.

\end{document}
